\scanpage{99}
que j'écris sous forme contragrédiente
\[
\check\kappa_i^{\otimes2}(\check\delta^{\otimes2})=
\check\nu(\check\delta^{\otimes2})\mathbin{\wedge}\omega_i\,,
\]
en termes des $\psi^\omega(\delta)$, comme
\[
\psi^\omega(\delta)^{\otimes2}=\check\nu(\check\delta^{\otimes2})\mathbin{\wedge}\omega\,. \tag*{(73)}
\]
\margin{\uline{NB} Moyennant (71) \inserted{(\unsure{resp.|ou} 72)}, les deux relations (73), pour les deux él. $\omega\in\underline\omega$, sont équivalentes, \unsure{resp.|il suffit} \deleted{\unsure{pour}} \inserted{sa} validité pour un $\delta$ \unsure{l'implique|l'impliquent} pour tous.}
On trouve ainsi le

\uline{Lemme}\quad Les épinglages de Legendre pour la donnée $(I,F_t\subset V_I(k),\underline t,\check\nu:\underline{\check t}^{\otimes2}\simeq F_t^{\otimes2}(\underline\omega))$, s'identifient aux systèmes $(\alpha,(\psi^\omega(\delta))_{\omega\in\underline\omega})$, où $\alpha:\mu_4^*\simeq\underline\omega$, et où $(\psi^\omega(\delta))_{\substack{\omega\in\underline\omega\\\delta\in\underline t^*}}\in(F_t)^{\underline\omega\times\underline t^*}$ \deleted{\unsure{satisf|satisfont|satisfaisant}} est une famille indexée par l'ens. \deleted{\unsure{à deux élé|des él.}} \inserted{produit $\underline\omega\times\underline t^*$} \unsure{ments|munts|notons} $\underline\omega=\underline\omega_I$, satisfaisant les conditions (71) (72) (73).

\deleted{\unsure{En}}Nous allons maintenant expliciter la donnée d'une famille à deux indices $(\psi_{\unsure{i|\delta}}^\omega(\delta))_{\substack{\omega\in\underline\omega\\\delta\in\underline t^\times}}$, à val. dans $F_t$, en termes du plongement
\[
F_t\subset V_I=\operatorname{Ker}(k^I\xrightarrow[\text{somme}]{}k)\,,
\]
on voit donc qu'elle s'exprime comme une famille à trois indices d'éléments de $k$ :
\[
\bigl(\psi_i^\omega(\delta)\bigr)_{\substack{i\in I\\\omega\in\underline\omega\\\delta\in\underline t^*}}\,,\qquad
\psi_i^\omega(\delta)\in k \tag*{(74)}
\]
Le fait que la famille d'éléments de $k^I$
% archive sheet number: 98
